\section{대칭 격자족: 이득 공식, 자유 작용의 한계, 비자유 작용}\label{sec:sym}

이 절은 3회차·5회차·7회차의 결과이다. 목표는 ``Klartag$\times$대칭'' 논증(AGZ, GGS)이 원리적으로 어디까지 갈 수 있는지를 군론·격자론의 언어로 정확히 서술하는 것이다.

\subsection{설정과 이득 항등식}
$G$ 를 유한군, $W$ 를 $\Q$-기약 $G$-가군, $A\subset\mathrm{End}_\Q(W)$ 를 $G$ 의 $\Q$-선형 생성(단순 대수), $\Delta=\mathrm{End}_G(W)$ (나눗셈대수)라 하자. 그러면 $A\cong M_r(\Delta^{op})$, $W\cong\Delta^r$, $\dim_\Q A=r^2\dim_\Q\Delta$, $\dim_\Q W=r\dim_\Q\Delta$ 이다. $G$ 가 $W^k$ ($n=k\dim_\Q W$)에 대각으로 작용할 때, $G$-불변 격자들은 $\Lambda=(W_\Z)^k\gamma$, $\gamma\in GL_k(\Delta_\R)$ 꼴의 족을 이루고 ($W_\Z\subset W$ 는 $G$-불변 격자, $\Delta_\R=\Delta\otimes\R$), $G$-불변 이차형식의 공간은 $\Delta_\R$ 위의 Hermite $k\times k$ 행렬 공간과 같다.

\begin{proposition}[이득 항등식]\label{prop:gain}\tagP\ (항등식) / \tagH\ (``이득'' 으로의 해석: 정리로 실현된 것은 자유 작용·$D_\R=\mathbb K^s$ 인 경우뿐, \S\ref{sec:quat})\
$\rho:=\chi_W/\chi_W(1)$ 이라 하면
\[
\sum_{g\in G}\rho(g)^2=\frac{|G|}{\dim_\Q A},\qquad
D_k:=\dim(\Sym^2 (W^k)_\R)^G=\frac1{|G|}\sum_{g}\frac{\chi(g)^2+\chi(g^2)}2=\frac{k^2\dim_\Q\Delta}{2}+O(k)\quad(\chi=k\chi_W).
\]
따라서 일반 궤도의 크기가 $w=|G|$ 이고 $-1\in G$ 일 때, 등변 Klartag 논증의 휴리스틱 부피 $wD_k/(2e)$ 의 Klartag ($w=2$, $D=N$) 대비 \emph{이득}은
\[
g(G,W)\;=\;\lim_{k\to\infty}\frac{wD_k}{2N}\;=\;\frac{|G|}{2\dim_\Q A}\;=\;\frac12\sum_{g\in G}\rho(g)^2 .
\]
\end{proposition}
\begin{proof}
$\chi_W$ 는 유리수 값이므로 $\sum_g\chi_W(g)^2=|G|\langle\chi_W,\chi_W\rangle=|G|\dim_\Q\mathrm{End}_G(W)=|G|\dim_\Q\Delta$. 이중 중심화 정리로 $\dim A\cdot\dim\Delta=(\dim W)^2$ 이므로 첫 식이 나온다. 둘째 식은 $\Sym^2$ 의 지표 공식이고, $\frac1{2|G|}\sum_g\chi(g)^2=k^2\dim\Delta/2$, 나머지 항은 $O(k)$ 이다. $N\sim k^2(\dim W)^2/2$ 를 대입하면 $wD_k/(2N)\to|G|\dim\Delta/(2(\dim W)^2)=|G|/(2\dim A)$.
\end{proof}
예: $\mu_m$ on $\Q(\zeta_m)$ ($m$ 짝수): $g=m/(2\varphi(m))$ (AGZ/GGS 의 $\log\log$). Hurwitz 단위군 $2T$ ($|G|=24$, $\dim A=4$): $g=3$. Icosian 단위군 $2I=SL_2(5)$ ($120$, $8$): $g=7.5$. $W$ 가 $\mathrm{Ind}_{\{\pm1\}}^G(\mathrm{sgn})$ 의 배수이면($-1$ 을 법으로 한 정칙표현; $\rho(\pm1)=\pm1$, 그 밖에서 0) $\tfrac12\sum\rho^2=1$: 이득 없음(기약이 아닌 경우로의 형식적 확장).

\subsection{고정점 없는 작용으로는 $\log\log$ 를 넘을 수 없다}
\begin{theorem}[자유 작용의 한계]\label{thm:fpf}\tagP\ (구조정리는 인용)\
$G$ 의 단위원 아닌 어떤 원소도 $W$ 에서 고윳값 1을 갖지 않는다고 하자. 그러면
\[
\sum_{g\in G}\rho(g)^2=\sum_{\substack{C\le G\ \text{순환}\\ |C|\ \text{제곱인수 없음}}}\frac1{\varphi(|C|)}\ \le\ 240\prod_{\ell\,\mid\,|G|}\frac \ell{\ell-1}\ =\ O(\log\log|G|)\qquad(\ell\text{ 는 소수}).
\]
즉 이득 지표는 $g\le120\prod_\ell\tfrac\ell{\ell-1}$ 이고, 자유 작용에 기반한 Venkatesh/AGZ 형 논증의 이득은 (상수배를 제외하고) 순환군이 이미 달성하는 $e^\gamma\log\log$ 를 넘지 못한다. (이는 회계량 $g$ 에 대한 정리이지 ``어떤 논증도 $\log\log$ 를 넘을 수 없다'' 는 장벽 정리가 아니다. 나눗셈대수의 유한부분군에 대한 같은 종류의 한계는 Gargava (2021)에 이미 있으며(1회차 문헌 노트), 여기서는 $A=M_r(\Delta)$ 인 자유 작용으로 확장했다.)
\end{theorem}
\begin{proof}
(i) 위수 $o$ 인 $g\ne1$ 의 고윳값 중 위수가 $o'<o$ 인 것이 있으면 $g^{o'}\neq1$ 이 고윳값 1을 가지므로 모순. 따라서 고윳값은 모두 원시 $o$ 제곱근이고, 유리성에 의해 $\chi_W(g)=\dim W\cdot\mu(o)/\varphi(o)$, $\rho(g)^2=\mu(o)^2/\varphi(o)^2$. 같은 순환부분군을 생성하는 원소가 $\varphi(o)$ 개이므로 등식이 나온다.

(ii) 이런 군은 Frobenius complement 의 구조를 가진다(위수 $\ell\ell'$ 인 부분군은 순환). Zassenhaus 의 구조정리(Passman, \emph{Permutation Groups}, \S18; 이 책은 이번 세션에서 열람하지 못했고 진술은 기억에 의존했으므로 확인이 필요하다): $G$ 는 모든 Sylow 부분군이 순환인 정규부분군 $M$ ($Z$-군)을 가지며 $[G:M]\le 240$ (가해인 경우 $G/M\hookrightarrow S_4$; 비가해인 경우 지수 $\le2$ 인 $SL_2(5)\times M$).

(iii) $Z$-군 $M=C_a\rtimes C_b$ ($\gcd(a,b)=1$)가 자유 작용하면 각 소수 위수의 부분군은 유일하다: $p\mid b$ 인 위수 $p$ 부분군 $P$ 와 $q\mid a$ 인 $Q\le C_a$ 에 대해 $PQ$ 는 위수 $pq$ 라서 순환, 즉 $P$ 는 $C_a$ 의 모든 소수 위수 부분군과 가환하고, 서로소 작용이므로 $C_a$ 전체와 가환한다. 따라서 $P$ 는 중심적이고 (Sylow 가 순환·공액이므로) 유일하다. 그러므로 $M$ 의 기약 성분 $W_1\subset W|_M$ 에 대해 $\sum_{m\in M}\rho_{W_1}(m)^2=\prod_{\ell\mid|M|}(1+\tfrac1{\ell-1})$.

(iv) $M$ 의 $\Q$-생성은 $A(G,W)$ 의 부분대수이고 $A(M,W_1)$ 로 전사하므로 $\dim A(G,W)\ge\dim A(M,W_1)$. 명제 \ref{prop:gain} 에 의해 $\sum_{G}\rho^2=|G|/\dim A(G,W)\le[G:M]\cdot|M|/\dim A(M,W_1)$.
\end{proof}
\begin{remark}
``자유 작용 $\Rightarrow$ $A$ 가 나눗셈대수'' 는 거짓이다 (3회차에서 발견된 정정): $SL_2(5)\times C_7$ 은 자유 작용하지만 $A=M_2(K)$, $K=\Q(\sqrt5,\zeta_7)$.
\end{remark}

\subsection{비자유 작용: 낮은 랭크 항과 텐서 안정성}
$r\ge2$ 이면 $\Lambda=W_\Z\otimes_{O_\Delta}L$ 은 모든 진부분 $\Delta$-부분공간 $U\subsetneq W$ 에 대해 부분격자 $\Lambda_U=(W_\Z\cap U)\otimes L$ (차원 $fn$, $f=\dim U/\dim W$)을 포함한다. 고정점을 갖는 원소의 고정공간 $W^g$ 가 그런 $U$ 의 예이다.

\begin{theorem}[낮은 랭크 항을 포함한 평균값 공식]\label{thm:meanvalue}\tagPL\ (Siegel--Weil--Rogers 공식 인용; $\Delta=K$ 수체, $W_\Z$ 자유인 경우만; 증명은 스케치 수준 + 수치 검산. 사원수 $\Delta$·class number 1 인 경우의 정확한 형태와 증명은 14·15회차: \S\ref{sec:leechthm})\
$k>r$ 이고 $L$ 이 랭크 $k$ 인 $O_K$-격자들의 Haar-랜덤 원소이면, 부피 $V\cdot\covol(\Lambda)$ 인 공 $B$ 에 대해
\[
\E\,\#\big((\Lambda\setminus0)\cap B\big)=V+\sum_{U\subsetneq W}T_U,\qquad
T_U=\frac{\Vol_{fn}(B\cap U_\R^k)}{c_U^{\,k}}=V^f\,\frac{(\pi n)^{f/2}}{(\pi fn)^{1/2}}\,R_U^{\,k/2}\,(1+O(\tfrac1{fn})),
\]
\[
R_U:=\Big(\frac{\dim W}{\dim U}\Big)^{\dim U}\frac{c_W^{\,2\dim U/\dim W}}{c_U^{\,2}},\qquad c_U:=\covol(W_\Z\cap U),\ c_W:=\covol(W_\Z).
\]
\end{theorem}
(증명: $X\ne0$ 의 열공간 $U$ 로 합을 분해하고 선형독립인 $s$-쌍에 대한 Siegel 공식을 각 $U$ 에 적용; \texttt{rounds/R3\_non\_free/notes.md} 정리 1, 명제 2. 독립 심사에서 수치 검산됨.)

\begin{definition}
$(G,W_\Z)$ 가 \emph{텐서 안정}: 모든 진부분 $\Delta$-부분공간 $U$ 에 대해 $R_U<1$. \emph{경계 안정}: $R_U\le1$.
동치로 $c_U^{2/\dim U}\dim U\ \ge\ c_W^{2/\dim W}\dim W$ (Grayson--Stuhler 반안정성보다 인자 $\dim W/\dim U$ 만큼 강한 조건).
\end{definition}
$R_U>1$ 인 $U$ 가 하나라도 있으면 $T_U$ 가 $k$ 에 대해 지수적으로 커져서 1차 모멘트 논증과 Klartag 형 논증이 모두 무너진다. $R_U=1$ (경계)이면 $T_U\asymp V^fn^{-(1-f)/2}$ 로 다항식적이다: $V\ll n$ 이면 무해하지만, $V\sim n^2$ 에서는 순수 1차 모멘트 논증이 더는 작동하지 않고($V=n^2$ 에서 $O(1)$ 이려면 $f\le1/5$), Klartag 형 과정에서 이 항이 접촉 예산 $wD\sim n^2$ 에 비해 무시할 만하다는 것은 \tagH\ 이다(증명 없음; 경계 부분공간의 개수 상계도 필요).

\textbf{주의(``불안정'' 의 의미).} ``텐서 불안정'' 은 1차 모멘트/Klartag 형 \emph{논증}이 그 족에서 작동하지 않는다는 뜻이며, 그 족에 조밀한 격자가 없다는 증명이 아니다(결정론적 모순에는 $R_U>4^{\dim U}$ 수준이 필요하다; 5회차 심사).

\begin{proposition}[No-go 정리들]\label{prop:nogo}\tagP
\begin{enumerate}
\item ($\Delta=\Q$, 텐서형.) $r=\dim W\ge2$ 이고 $\Delta=\Q$ 이면 \emph{텐서형} 족 $W_\Z\otimes_\Z L$ 은 어떤 $W_\Z$ 에 대해서도 최소벡터 방향 직선 $U$ 에서 $R_U^{-1}=\gamma(W_\Z)/r\le 2/\pi<1$, 즉 항상 불안정. (심사에서 범위가 정정됨: 텐서형이 아닌 $G$-불변 격자 --- 비극대 order $\Z[G]$ 의 가군 --- 는 이 논법으로 배제되지 않는다. 예: $D_8$ 에 대해 $\{(x,y)\in L_1^2: y\equiv\phi(x)\bmod 2L_1\}$ 는 모든 직선에서 $R_U\le1$. 다만 이런 족은 자유 $\mu_4$/$\mu_6$ 족의 부분족이라 이득이 자유 작용의 값과 같다. 반사나 $a\ne b$ 인 대합을 포함하는 군(Weyl 군, $S_n$, 실 Clifford 군)은 아래 (3)으로 \emph{모든} 불변 격자에 대해 배제된다.)
\item (판별식.) 직선 조건은 $\gamma(W_\Z)\ge\dim_\Q W/(t\cdot\rd_\Delta)$ 를 함의한다 ($t^2=\dim_{Z(\Delta)}\Delta$, $\rd_\Delta$ 는 극대 order 의 정규화된 root discriminant). 큰 차원에서는 $t\,\rd_\Delta\ge 2\pi e/1.744\approx 9.8$ 이 필요하다.
\item (엔트로피 부등식.) 소수 위수 $q$ 의 $g\in G$ 가 $W|_{\langle g\rangle}\cong\Q^a\oplus\Q(\zeta_q)^b$ ($a>0$)이면, 안정성은
\[(a+b(q-1))\,h\!\Big(\frac a{a+b(q-1)}\Big)\ \le\ 2\min(a,b)\log q\qquad(h(x)=-x\log x-(1-x)\log(1-x))\]
를 함의한다. $q=2$ 이면 $a=b$ 이고 정확히 경계이다(중심이 아닌 대합은 결코 엄격 안정일 수 없다).
\item (결정론적 손실.) $R_U^{1/\dim U}>2^{0.802}$ 인 $U$ 가 있으면 \emph{어떤} $L$ 에 대해서도 $V(W_\Z\otimes L)$ 은 Minkowski--Hlawka 보다 지수적으로 작다.
\end{enumerate}
\end{proposition}
(증명은 \texttt{rounds/R3\_non\_free/notes.md} 명제 3--5, 따름정리 4.1. (3)의 핵심: $I:=[W_\Z:(W_\Z\cap U)\oplus(W_\Z\cap U')]\le q^{\min(a,b)}$ 이고 $c_Uc_{U'}=I\,c_W$; $U$ 와 $U'=\ker\Phi_q(g)$ 에 대한 안정성 부등식을 곱한다.)

이번 세션에서 추가로 얻은, unimodular 격자에 대한 더 강한 형태(저자의 유도와 5회차 이론 에이전트의 Lemma A$'$ 가 독립적으로 일치했고 5회차 심사를 통과):
\begin{lemma}[Unimodular 강직성]\label{lem:unimod}\tagP\
$W_\Z$ 가 ($G$-불변 내적에 대해) unimodular 이고, $g\in G$ 가 소수 위수 $q$, $W|_{\langle g\rangle}\cong\Q^a\oplus\Q(\zeta_q)^b$, $a>0$ 이라 하자. 그러면 $U=W^g$ 에 대해 $R_U\ge1$ 이고, 등호는 정확히 $a=b$ 이고 접착 지수가 최대값 $q^a$ 일 때(즉 $W_\Z$ 가 $\Z[C_q]$-사영(국소 자유)가군일 때) 성립한다. 즉 unimodular 인 경우 고정점을 갖는 소수 위수 원소는 ``정칙표현의 배수 + 최대 접착'' 으로 작용할 때에만 허용되고, 그때 정확히 경계이다.
\end{lemma}
\begin{proof}
$W_\Z$ 가 unimodular 이므로 $U$ 로의 직교사영 $e_gW_\Z$ 는 $W_\Z\cap U$ 의 ($U$ 안에서의) 쌍대격자이고, $\det(W_\Z\cap U)=[e_gW_\Z:W_\Z\cap U]$. $q\,e_g=\sum_jg^j\in\Z[g]$ 이므로 이 지수는 $q^{a}$ 를 나누고, 같은 지수가 $[e'_gW_\Z:W_\Z\cap U']$ 이기도 한데 $(1-g)$ 가 이 몫을 죽이므로 $q^{b}$ 도 나눈다: $\det(W_\Z\cap U)\le q^{\min(a,b)}$. 한편 $R_U\le1$ 은 $\det(W_\Z\cap U)\ge\big(\tfrac{a+b(q-1)}a\big)^a$ 와 동치. $a\le b$ 이면 $q\ge1+(q-1)b/a$, 즉 $a\ge b$ 가 필요하므로 $a=b$ 이고 지수가 최대($q^a$, 즉 $\Z[C_q]$-사영). $a>b$ 이면 $t=a/b>1$ 에 대해 $\log q\ge t\log(1+(q-1)/t)$ 가 필요한데, 우변은 $t$ 에 대해 순증가하고 $t=1$ 에서 $\log q$ 이므로 불가능.
\end{proof}
이 보조정리는 ``경계 안정성'' 이 왜 모든 예에서 정확히 $R_U=1$ 로 나타나는지 설명해 준다 (아래 표, 그리고 \S\ref{sec:leech}).

\subsection{예}
\begin{center}
\begin{tabular}{@{}llll@{}}
\toprule
$W_\Z$, $G$, $\Delta$, $r$ & $U$ & $f\tau_U^2:=R_U^{-1/\dim U}$ & 판정 \\
\midrule
$\Z^2$, $D_8$, $\Q$, 2 & 직선 $(1,0)$ & 0.500 & 실패 \\
$A_2$, $D_{12}$, $\Q$, 2 & 최소 직선 & 0.577 & 실패 \\
$D_4$, $2^{1+4}_+$, $\Q$, 4 & 최소 직선 & 0.354 & 실패 \\
$E_8$, $2^{1+6}_+$, $\Q$, 8 & 최소 직선 & 0.250 & 실패 \\
$O_K^p$, Heisenberg $p^{1+2}$, $\Q(\zeta_p)$, $p$ & $Ke_1$ & $1/p$ & 실패 \\
Craig $A_{p-1}^{((p+1)/4)}$, $PSL_2(p)$, $\Q(\sqrt{-p})$ & 최소 직선 & $\to 1/2$ & 실패 \\
$D_4$ ($\Z[\sqrt{-2}]$-격자, 랭크 2) & 최소 $K$-직선 & 1.000 & 경계 \\
$O_E/O_K$ 비분기 확대 ($E=\Q(\sqrt{-3},\sqrt5)$, $K=\Q(\sqrt{-15})$) & $U=K$ & 1.000 & 경계 \\
Leech, $SL_2(13)$, $(-1,-1)_{\Q(\sqrt{13})}$, 3 & 모든 직선·평면 & $\ge 1.000$ & \textbf{경계} (\S\ref{sec:leech}) \\
\bottomrule
\end{tabular}
\end{center}
($f\tau_U^2\ge1\iff R_U\le1$. 수치는 \texttt{rounds/R3\_non\_free/check\_nonfree.py} 및 \texttt{rounds/R5\_leech13/} 의 출력.)
